\documentclass[11pt]{letter}
\usepackage[margin=1in]{geometry}
\usepackage{hyperref}
\usepackage{url}

\signature{Jeffrey D.\ Varner, Ph.D.\\
Robert Frederick Smith School of Chemical and Biomolecular Engineering\\
Cornell University\\
Ithaca, NY 14853\\
\texttt{jdv27@cornell.edu}}

\date{\today}

\begin{document}

\begin{letter}{Editorial Board\\
\textit{npj Systems Biology and Applications}}

\opening{Dear Editors,}

We are pleased to submit our manuscript entitled
\textbf{``Validated Synthetic Patient Generation from Small Longitudinal
Cohorts Using Stochastic Attention''}
for consideration as an Article in \textit{npj Systems Biology and Applications}.

Small longitudinal clinical cohorts, common in maternal health and rare disease
research, present a fundamental challenge: too few patients to support
computational modeling, yet too costly and slow to expand through additional
enrollment. In this work, we address this gap with multiplicity-weighted
Stochastic Attention (SA), a generative framework rooted in modern Hopfield
network theory that produces synthetic patients by interpolating between
stored patient profiles on a continuous energy landscape. Unlike parametric
methods such as multivariate normal sampling or deep generative models (GANs,
VAEs), SA operates directly on the geometry of the observed cohort and requires
no parameter estimation, making it well suited to the $n < p$ regime that
characterizes most small clinical studies.

We applied SA to a longitudinal coagulation dataset from 23 pregnant patients
spanning 72 biochemical features across three visits, including rare subgroups
(polycystic ovary syndrome, $n{=}3$; preeclampsia, $n{=}5$). We validated the
synthetic patients through four progressively stringent levels: marginal
plausibility, cross-visit covariance preservation, conditional subgroup
amplification, and mechanistic consistency with a separately specified,
generator-blind 58-ODE model calibrated on the same cohort. Across these
validation levels, the synthetic patients closely matched their real counterparts
under the reported metrics; non-significant difference tests are treated as
descriptive non-detection rather than proof of equivalence. A downstream utility
test further showed that a mechanistic model calibrated entirely on synthetic
patients predicted held-out real patient outcomes as well as one calibrated on
real data.

We believe this work is well suited to \textit{npj Systems Biology and
Applications} for several reasons. First, the methodology bridges systems
biology (mechanistic ODE modeling of the coagulation cascade) with modern
machine learning (Hopfield networks, Langevin dynamics), a combination that
aligns with the journal's scope. Second, the multi-level validation framework,
particularly the mechanistic validation using a separately specified,
generator-blind ODE model calibrated on the same cohort, provides
a template that transfers to any domain with a calibrated mechanistic model
(pharmacokinetic, physiological, metabolic). Third, the multiplicity-weighted
conditional generation capability addresses a practical need in rare disease and
maternal health research, where generating condition-specific synthetic cohorts
from very small subgroups supports exploratory simulation and power-analysis
workflows without increasing the subgroup's effective inferential sample size or
replacing recruitment.

A preprint of this work is available on arXiv (arXiv:2604.07557). The
manuscript has not been submitted elsewhere.

All code, validation scripts, and generated datasets are publicly available at
\url{https://github.com/varnerlab/SA-generation-legacy-dataset-paper}.

We confirm that all authors have approved the manuscript and that there are no
competing interests to declare.

Thank you for considering our submission. We look forward to your response.

\closing{Sincerely,}

\end{letter}
\end{document}
